\section{Introduction}\label{sec:introduction}

Mean--variance allocation requires a covariance matrix, yet that
matrix is hardest to estimate in the settings where diversification
matters most.
An unrestricted covariance matrix for $n$ assets contains $n(n+1)/2$
entries, while a demeaned return history of length $T$ has rank at
most $\min(T-1,n)$.
Short histories therefore leave a portfolio manager with two linked
problems: the strongest sample directions are noisy, and the
optimizer is most sensitive to the weakest ones.
Shrinkage and factor models reduce this burden by imposing structure
\citep{ledoit_wolf_2004,kelly_characteristics_2019}, but their
cross-asset information still originates in joint returns.

An alternative is to ask how much portfolio risk observable firm
information can rule out before a return covariance matrix is estimated.
We represent each firm by the distribution of its article embeddings
and use quadratic optimal transport to compare those distributions.
When two information distributions are sufficiently different,
maintained restrictions linking information to systematic exposures
imply that the firms' latent risks cannot be perfectly aligned.
The resulting separation certifies a diversification benefit relative
to perfect positive dependence.

For a standardized long-only portfolio with normalized risk weights
$q$, the headline result takes the form
\[
  \operatorname{Var}(R_q)\le 1-\mathcal C(q).
\]
The value $1$ is the variance benchmark under perfect positive dependence,
and $\mathcal C(q)$ is the amount that the observed information
geometry certifies away.
A larger certificate therefore tightens the admissible upper bound on
portfolio risk.
It does not estimate the covariance matrix.

\begin{figure}
  \centering
  \ppfigure{0.95}{conceptual_pipeline}
  \caption{Observed information to portfolio choice. Solid arrows connect
    computed objects; the dashed arrow from observed W2 separation to the
    latent floor is the maintained transmission restriction. Blue denotes
    observed objects, orange latent and certificate objects, green the risk
  bound, and purple the allocation decision.}
  \label{fig:p3-conceptual-pipeline}
\end{figure}

The same distribution-valued representation has generated related
financial objects at neighbouring levels of aggregation.
At the pairwise level, \citet{gawronsky_continuous_2026} derive a
covariance envelope from Wasserstein separation.
At the cross-sectional level, \citet{gawronsky_spatial_2027} use
target-anchored Wasserstein barycentric reconstruction to form a
barycentric interaction field.
The contribution here is at the portfolio level: it aggregates
information-derived separation within one coherent joint exposure law
and turns the resulting variance bound into a decision rule.
The construction below is self-contained.

Let $C_i$ denote firm $i$'s observed characteristic law and let $P_i$
denote its latent exposure law in factor-risk coordinates.
Empirically, $C_i$ is the distribution of the firm's article embeddings.
The quadratic Wasserstein distance $W_2(C_i,C_j)$ is the minimum
root-mean-square displacement required to match the two article clouds.
It therefore measures how much semantic mass must be rearranged to
make the observed information distributions coincide.
The latent law $P_i$ describes systematic exposure rather than text,
so $C_i$ and $P_i$ need not share units and one does not identify the
other directly.

Three maintained links carry the argument from observed information
to portfolio risk.
A common information-to-exposure map prevents distinct information
states from collapsing into identical systematic exposures, while
firm-specific slack permits bounded departures from that common map.
A coherent joint law ensures that all pairwise exposure relations can
coexist within the same portfolio.
A return bridge then connects systematic exposure variance to
standardized total returns.
The text determines the observed geometry; it does not identify these
transmission restrictions.

For a common carrier constant $L>0$ and firm-specific slack radii
$\tau_i$, the observable pairwise floor is
\[
  \ell_{ij}
  =\left[L^{-1}W_2(C_i,C_j)-\tau_i-\tau_j\right]_+.
\]
This floor is the exposure separation that remains after allowing for
common distortion and the two firms' deviations from the common map.
The positive part records that the observed distance may be too small
to certify any separation once slack is deducted.
The portfolio certificate aggregates these floors using normalized
risk weights for a long-only portfolio:
\[
  \mathcal C(q)
  =\frac12\sum_i \sum_j q_i q_j\ell_{ij}^2.
\]
Separation between a pair contributes only when the portfolio holds
both firms, and contributes more when their joint portfolio weight is larger.
If the exposure marginals belong to one coherent joint law and the
maintained carrier and slack restrictions hold, systematic portfolio
variance is bounded by weighted marginal second moments less $\mathcal C(q)$.
The standardized return bridge then gives the headline bound above.
Observed differences between firms therefore restrict the joint risk
configurations that remain admissible.

Investors choose capital weights rather than normalized risk weights.
For capital weights $x_i$ and marginal volatility scales $\sigma_i$,
define $A(x)=\sum_i x_i\sigma_i$ and $q_i(x)=x_i\sigma_i/A(x)$.
The corresponding raw-return certificate is
\[
  \operatorname{Var}(R_x)
  \le A(x)^2\{1-\mathcal C(q(x))\}.
\]
Minimizing this upper bound yields an information-certified portfolio without
an expected-return input.
Compactness of the feasible long-only set supplies existence of a
minimizer, while a directly checkable condition on the observed
distance matrix makes the standardized objective convex.

The empirical exercise keeps construction separate from evaluation.
The canonical zero-slack implementation, which we call the news-only
allocation, is formed from observed W2 geometry without using cross-asset return
covariance and is evaluated only afterwards on the full-sample standardized
covariance matrix.
Across four prespecified capped long-only reference populations,
between \ppnum[2]{news-only-percentile-uniform-cap-15-pct}\% and
\ppnum[2]{news-only-percentile-concentrated-cap-15-pct}\% of feasible
portfolios have variance no greater than the news-only allocation.
The corresponding share for equal risk weights---the inverse-volatility capital
benchmark expressed in normalized risk-weight coordinates---lies between
\ppnum[1]{equal-percentile-concentrated-cap-15-pct}\% and
\ppnum[1]{equal-percentile-uniform-cap-12p5-pct}\%.
These rankings are descriptive and in-sample; they illustrate the
allocation implied by the maintained model rather than forecast
out-of-sample performance.

The paper makes three contributions.
First, it derives a coherent portfolio variance bound in a sharp
multi-firm form and supplies the computationally simpler weighted
pairwise relaxation used by the decision rule.
Second, it shows that when firm-specific slack is zero, the common
carrier scale changes the certified variance reduction but not the
normalized allocation, which is determined by the observed information geometry.
Third, in a \ppnum[0]{n-tickers}-firm in-sample exercise, it reports
how the allocation's
conventionally evaluated variance ranks under four prespecified
feasible-portfolio laws.

The argument proceeds from identification to decision and then to
descriptive evaluation.
\Cref{sec:literature} positions the contribution in structured
covariance, factor-risk, and textual-characteristic research.
\Cref{sec:model} defines the observed, latent, and coherent portfolio
objects and derives the information-certified variance bound.
\Cref{sec:portfolio-choice} turns the pairwise certificate into a decision rule.
\Cref{sec:empirical-design,sec:results} describe the data and report
the in-sample variance ranking.
\Cref{sec:limitations} discusses what the results establish and where
their boundaries lie before \Cref{sec:conclusion} concludes.
